\documentclass[11pt]{article}
\usepackage[margin=0.84in]{geometry}
\usepackage{amsmath,amssymb,mathtools,booktabs,tabularx,microtype,fancyhdr,enumitem}
\usepackage[dvipsnames]{xcolor}\usepackage[colorlinks=true,linkcolor=NavyBlue,urlcolor=BrickRed,hypertexnames=false]{hyperref}
\definecolor{ink}{HTML}{172337}\definecolor{accent}{HTML}{8B4A35}\definecolor{pale}{HTML}{FAF1ED}\definecolor{passgreen}{HTML}{147A4A}
\color{ink}\pagestyle{fancy}\fancyhf{}\lhead{$PSL_2(q)$ on the projective line}\rhead{Proof and verification}\cfoot{\thepage}\setlength{\headheight}{14pt}
\newcommand{\Sym}{\operatorname{Sym}}\newcommand{\Exp}{\operatorname{Exp}_{\sigma}}\newcommand{\PASS}{\textcolor{passgreen}{\textbf{PASS}}}
\newcommand{\summarybox}[1]{\begin{center}\fcolorbox{accent}{pale}{\parbox{0.91\linewidth}{#1}}\end{center}}
\title{\textbf{$PSL_2(q)$ on $\mathbf P^1(\mathbf F_q)$}\\[3pt]\large Exact sigma--Molien formula, transitivity boundary, and verification}
\author{Computational verification note}\date{17 July 2026}
\begin{document}\maketitle
\summarybox{\textbf{Outcome.} Formula (3.3) in the supplied statement is correct.  Prime-field groups $q=2,3,5,7$ were constructed from all determinant-one $2\times2$ matrices and then quotiented by their projective action.  All cycle-type multiplicities, $48$ coefficients through degree four, and direct orbit counts agree.  Determinant-level error is at most $6.67\times10^{-16}$.}

\section{Representation and coefficient meaning}
Let $q=p^f$, $d=\gcd(2,q-1)$, and $G=PSL_2(q)$, so $|G|=q(q^2-1)/d$.  Let $V=\mathbb C[\mathbf P^1(\mathbf F_q)]$, of dimension $q+1$.  For a permutation with $c_m$ cycles of length $m$,
\[
 \det(I-tg)^{-1}=\prod_m(1-t^m)^{-c_m},\qquad
 \prod_j\det(I-t_jg)^{-1}=\Exp\left(\sum_m c_mp_m\right). \tag{1}
\]
The coefficient at $\tau$ equals
$\dim(\bigotimes_s\Sym^{\tau_s}V)^G$ by the Reynolds-projector argument.  Because $V$ has a basis permuted by $G$, it also equals the number of $G$-orbits on tuples of multisets having sizes $\tau_s$.  This gives a genuinely independent combinatorial check.

\section{Four element types}
Put $L_-=(q-1)/d$ and $L_+=(q+1)/d$.
\begin{center}\begin{tabularx}{\linewidth}{@{}lXl@{}}\toprule
type & cycle structure on $\mathbf P^1$ & multiplicity\\\midrule
identity & $1^{q+1}$ & $1$\\
nonidentity unipotent & $1^1p^{q/p}$ & $q^2-1$\\
split semisimple, order $m>1$ & $1^2m^{(q-1)/m}$ & $\frac{q(q+1)}2\varphi(m)$, $m\mid L_-$\\
nonsplit semisimple, order $m>1$ & $m^{(q+1)/m}$ & $\frac{q(q-1)}2\varphi(m)$, $m\mid L_+$\\\bottomrule
\end{tabularx}\end{center}
The unipotent form $x\mapsto x+1$ fixes infinity and partitions the affine line into $q/p$ cycles.  A split element has two eigenlines; a nonsplit element has none.  Each nontrivial semisimple element determines its eigenline pair, hence its unique torus.  The split and nonsplit torus counts are $q(q+1)/2$ and $q(q-1)/2$ respectively, giving the displayed multiplicities.

\section{Exact formula}
Substituting the four rows into (1) proves
\begin{align}
\mathfrak M_{G,V}=\frac1{|G|}\Bigg[&\Exp((q+1)p_1)
+(q^2-1)\Exp\left(p_1+\frac qp p_p\right)\notag\\
&+\frac{q(q+1)}2\sum_{\substack{m\mid L_-\\m>1}}\varphi(m)
 \Exp\left(2p_1+\frac{q-1}{m}p_m\right)\notag\\
&+\frac{q(q-1)}2\sum_{\substack{m\mid L_+\\m>1}}\varphi(m)
 \Exp\left(\frac{q+1}{m}p_m\right)\Bigg]. \tag{2}
\end{align}
The multiplicities sum to $|G|$, an immediate global consistency check.

\section{Poisson comparison and the exact boundary}
Two-transitivity forces the same orbit counts as the stable symmetric-group action through total degree two:
\[
 [m_{(1)}]=1,\qquad [m_{(2)}]=[m_{(1,1)}]=2.
\]
For ordered triples, equality patterns contribute $1+3$ orbits, while distinct triples contribute
\[
 \frac{(q+1)q(q-1)}{|G|}=d.
\]
Thus
\begin{equation}[m_{(1,1,1)}]=4+d=\begin{cases}5&q\text{ even},\\6&q\text{ odd}.\end{cases} \tag{3}\end{equation}
At degree four, equality patterns with $1,2,3,4$ distinct values give
\begin{equation}[m_{(1,1,1,1)}]=1+7+6d+d(q-2)=8+d(q+4). \tag{4}\end{equation}
The last term is the projective cross-ratio parameter.  Consequently the raw sigma--MGF has no finite coefficientwise limit as $q\to\infty$.

\section{Verification strategy}
For each prime $q$ tested, the implementation enumerates every
$(a,b,c,d)\in\mathbf F_q^4$ with $ad-bc=1$, induces the fractional-linear
permutation on $\mathbf F_q\cup\{\infty\}$, and removes scalar-kernel
duplicates.  It compares the actual cycle-type histogram with the four-type
theorem, extracts every coefficient through degree four, independently counts
orbits on tuples of multisets, and finally compares determinant averages of
the permutation matrices at $(t_1,t_2)=(0.07,0.11)$.
The construction is intentionally transparent and is presently implemented for prime fields.  The proof and formula cover all prime powers; testing $f>1$ would require a small finite-field layer, not a change to the group-theoretic argument.

\end{document}
